
%% bare_jrnl_compsoc.tex
%% V1.4b
%% 2015/08/26
%% by Michael Shell
%% See:
%% http://www.michaelshell.org/
%% for current contact information.
%%
%% This is a skeleton file demonstrating the use of IEEEtran.cls
%% (requires IEEEtran.cls version 1.8b or later) with an IEEE
%% Computer Society journal paper.
%%
%% Support sites:
%% http://www.michaelshell.org/tex/ieeetran/
%% http://www.ctan.org/pkg/ieeetran
%% and
%% http://www.ieee.org/

%%*************************************************************************
%% Legal Notice:
%% This code is offered as-is without any warranty either expressed or
%% implied; without even the implied warranty of MERCHANTABILITY or
%% FITNESS FOR A PARTICULAR PURPOSE! 
%% User assumes all risk.
%% In no event shall the IEEE or any contributor to this code be liable for
%% any damages or losses, including, but not limited to, incidental,
%% consequential, or any other damages, resulting from the use or misuse
%% of any information contained here.
%%
%% All comments are the opinions of their respective authors and are not
%% necessarily endorsed by the IEEE.
%%
%% This work is distributed under the LaTeX Project Public License (LPPL)
%% ( http://www.latex-project.org/ ) version 1.3, and may be freely used,
%% distributed and modified. A copy of the LPPL, version 1.3, is included
%% in the base LaTeX documentation of all distributions of LaTeX released
%% 2003/12/01 or later.
%% Retain all contribution notices and credits.
%% ** Modified files should be clearly indicated as such, including  **
%% ** renaming them and changing author support contact information. **
%%*************************************************************************


% *** Authors should verify (and, if needed, correct) their LaTeX system  ***
% *** with the testflow diagnostic prior to trusting their LaTeX platform ***
% *** with production work. The IEEE's font choices and paper sizes can   ***
% *** trigger bugs that do not appear when using other class files.       ***                          ***
% The testflow support page is at:
% http://www.michaelshell.org/tex/testflow/


\documentclass[10pt,journal,compsoc]{IEEEtran}
%
% If IEEEtran.cls has not been installed into the LaTeX system files,
% manually specify the path to it like:
% \documentclass[10pt,journal,compsoc]{../sty/IEEEtran}





% Some very useful LaTeX packages include:
% (uncomment the ones you want to load)


% *** MISC UTILITY PACKAGES ***
%
%\usepackage{ifpdf}
% Heiko Oberdiek's ifpdf.sty is very useful if you need conditional
% compilation based on whether the output is pdf or dvi.
% usage:
% \ifpdf
%   % pdf code
% \else
%   % dvi code
% \fi
% The latest version of ifpdf.sty can be obtained from:
% http://www.ctan.org/pkg/ifpdf
% Also, note that IEEEtran.cls V1.7 and later provides a builtin
% \ifCLASSINFOpdf conditional that works the same way.
% When switching from latex to pdflatex and vice-versa, the compiler may
% have to be run twice to clear warning/error messages.






% *** CITATION PACKAGES ***
%
\ifCLASSOPTIONcompsoc
  % IEEE Computer Society needs nocompress option
  % requires cite.sty v4.0 or later (November 2003)
  \usepackage[nocompress]{cite}
\else
  % normal IEEE
  \usepackage{cite}
\fi
% cite.sty was written by Donald Arseneau
% V1.6 and later of IEEEtran pre-defines the format of the cite.sty package
% \cite{} output to follow that of the IEEE. Loading the cite package will
% result in citation numbers being automatically sorted and properly
% "compressed/ranged". e.g., [1], [9], [2], [7], [5], [6] without using
% cite.sty will become [1], [2], [5]--[7], [9] using cite.sty. cite.sty's
% \cite will automatically add leading space, if needed. Use cite.sty's
% noadjust option (cite.sty V3.8 and later) if you want to turn this off
% such as if a citation ever needs to be enclosed in parenthesis.
% cite.sty is already installed on most LaTeX systems. Be sure and use
% version 5.0 (2009-03-20) and later if using hyperref.sty.
% The latest version can be obtained at:
% http://www.ctan.org/pkg/cite
% The documentation is contained in the cite.sty file itself.
%
% Note that some packages require special options to format as the Computer
% Society requires. In particular, Computer Society  papers do not use
% compressed citation ranges as is done in typical IEEE papers
% (e.g., [1]-[4]). Instead, they list every citation separately in order
% (e.g., [1], [2], [3], [4]). To get the latter we need to load the cite
% package with the nocompress option which is supported by cite.sty v4.0
% and later. Note also the use of a CLASSOPTION conditional provided by
% IEEEtran.cls V1.7 and later.





% *** GRAPHICS RELATED PACKAGES ***
%
\ifCLASSINFOpdf
  % \usepackage[pdftex]{graphicx}
  % declare the path(s) where your graphic files are
  % \graphicspath{{../pdf/}{../jpeg/}}
  % and their extensions so you won't have to specify these with
  % every instance of \includegraphics
  % \DeclareGraphicsExtensions{.pdf,.jpeg,.png}
\else
  % or other class option (dvipsone, dvipdf, if not using dvips). graphicx
  % will default to the driver specified in the system graphics.cfg if no
  % driver is specified.
  % \usepackage[dvips]{graphicx}
  % declare the path(s) where your graphic files are
  % \graphicspath{{../eps/}}
  % and their extensions so you won't have to specify these with
  % every instance of \includegraphics
  % \DeclareGraphicsExtensions{.eps}
\fi
% graphicx was written by David Carlisle and Sebastian Rahtz. It is
% required if you want graphics, photos, etc. graphicx.sty is already
% installed on most LaTeX systems. The latest version and documentation
% can be obtained at: 
% http://www.ctan.org/pkg/graphicx
% Another good source of documentation is "Using Imported Graphics in
% LaTeX2e" by Keith Reckdahl which can be found at:
% http://www.ctan.org/pkg/epslatex
%
% latex, and pdflatex in dvi mode, support graphics in encapsulated
% postscript (.eps) format. pdflatex in pdf mode supports graphics
% in .pdf, .jpeg, .png and .mps (metapost) formats. Users should ensure
% that all non-photo figures use a vector format (.eps, .pdf, .mps) and
% not a bitmapped formats (.jpeg, .png). The IEEE frowns on bitmapped formats
% which can result in "jaggedy"/blurry rendering of lines and letters as
% well as large increases in file sizes.
%
% You can find documentation about the pdfTeX application at:
% http://www.tug.org/applications/pdftex






% *** MATH PACKAGES ***
%
%\usepackage{amsmath}
% A popular package from the American Mathematical Society that provides
% many useful and powerful commands for dealing with mathematics.
%
% Note that the amsmath package sets \interdisplaylinepenalty to 10000
% thus preventing page breaks from occurring within multiline equations. Use:
%\interdisplaylinepenalty=2500
% after loading amsmath to restore such page breaks as IEEEtran.cls normally
% does. amsmath.sty is already installed on most LaTeX systems. The latest
% version and documentation can be obtained at:
% http://www.ctan.org/pkg/amsmath





% *** SPECIALIZED LIST PACKAGES ***
%
%\usepackage{algorithmic}
% algorithmic.sty was written by Peter Williams and Rogerio Brito.
% This package provides an algorithmic environment fo describing algorithms.
% You can use the algorithmic environment in-text or within a figure
% environment to provide for a floating algorithm. Do NOT use the algorithm
% floating environment provided by algorithm.sty (by the same authors) or
% algorithm2e.sty (by Christophe Fiorio) as the IEEE does not use dedicated
% algorithm float types and packages that provide these will not provide
% correct IEEE style captions. The latest version and documentation of
% algorithmic.sty can be obtained at:
% http://www.ctan.org/pkg/algorithms
% Also of interest may be the (relatively newer and more customizable)
% algorithmicx.sty package by Szasz Janos:
% http://www.ctan.org/pkg/algorithmicx




% *** ALIGNMENT PACKAGES ***
%
%\usepackage{array}
% Frank Mittelbach's and David Carlisle's array.sty patches and improves
% the standard LaTeX2e array and tabular environments to provide better
% appearance and additional user controls. As the default LaTeX2e table
% generation code is lacking to the point of almost being broken with
% respect to the quality of the end results, all users are strongly
% advised to use an enhanced (at the very least that provided by array.sty)
% set of table tools. array.sty is already installed on most systems. The
% latest version and documentation can be obtained at:
% http://www.ctan.org/pkg/array


% IEEEtran contains the IEEEeqnarray family of commands that can be used to
% generate multiline equations as well as matrices, tables, etc., of high
% quality.




% *** SUBFIGURE PACKAGES ***
%\ifCLASSOPTIONcompsoc
%  \usepackage[caption=false,font=footnotesize,labelfont=sf,textfont=sf]{subfig}
%\else
%  \usepackage[caption=false,font=footnotesize]{subfig}
%\fi
% subfig.sty, written by Steven Douglas Cochran, is the modern replacement
% for subfigure.sty, the latter of which is no longer maintained and is
% incompatible with some LaTeX packages including fixltx2e. However,
% subfig.sty requires and automatically loads Axel Sommerfeldt's caption.sty
% which will override IEEEtran.cls' handling of captions and this will result
% in non-IEEE style figure/table captions. To prevent this problem, be sure
% and invoke subfig.sty's "caption=false" package option (available since
% subfig.sty version 1.3, 2005/06/28) as this is will preserve IEEEtran.cls
% handling of captions.
% Note that the Computer Society format requires a sans serif font rather
% than the serif font used in traditional IEEE formatting and thus the need
% to invoke different subfig.sty package options depending on whether
% compsoc mode has been enabled.
%
% The latest version and documentation of subfig.sty can be obtained at:
% http://www.ctan.org/pkg/subfig




% *** FLOAT PACKAGES ***
%
%\usepackage{fixltx2e}
% fixltx2e, the successor to the earlier fix2col.sty, was written by
% Frank Mittelbach and David Carlisle. This package corrects a few problems
% in the LaTeX2e kernel, the most notable of which is that in current
% LaTeX2e releases, the ordering of single and double column floats is not
% guaranteed to be preserved. Thus, an unpatched LaTeX2e can allow a
% single column figure to be placed prior to an earlier double column
% figure.
% Be aware that LaTeX2e kernels dated 2015 and later have fixltx2e.sty's
% corrections already built into the system in which case a warning will
% be issued if an attempt is made to load fixltx2e.sty as it is no longer
% needed.
% The latest version and documentation can be found at:
% http://www.ctan.org/pkg/fixltx2e


%\usepackage{stfloats}
% stfloats.sty was written by Sigitas Tolusis. This package gives LaTeX2e
% the ability to do double column floats at the bottom of the page as well
% as the top. (e.g., "\begin{figure*}[!b]" is not normally possible in
% LaTeX2e). It also provides a command:
%\fnbelowfloat
% to enable the placement of footnotes below bottom floats (the standard
% LaTeX2e kernel puts them above bottom floats). This is an invasive package
% which rewrites many portions of the LaTeX2e float routines. It may not work
% with other packages that modify the LaTeX2e float routines. The latest
% version and documentation can be obtained at:
% http://www.ctan.org/pkg/stfloats
% Do not use the stfloats baselinefloat ability as the IEEE does not allow
% \baselineskip to stretch. Authors submitting work to the IEEE should note
% that the IEEE rarely uses double column equations and that authors should try
% to avoid such use. Do not be tempted to use the cuted.sty or midfloat.sty
% packages (also by Sigitas Tolusis) as the IEEE does not format its papers in
% such ways.
% Do not attempt to use stfloats with fixltx2e as they are incompatible.
% Instead, use Morten Hogholm'a dblfloatfix which combines the features
% of both fixltx2e and stfloats:
%
% \usepackage{dblfloatfix}
% The latest version can be found at:
% http://www.ctan.org/pkg/dblfloatfix




%\ifCLASSOPTIONcaptionsoff
%  \usepackage[nomarkers]{endfloat}
% \let\MYoriglatexcaption\caption
% \renewcommand{\caption}[2][\relax]{\MYoriglatexcaption[#2]{#2}}
%\fi
% endfloat.sty was written by James Darrell McCauley, Jeff Goldberg and 
% Axel Sommerfeldt. This package may be useful when used in conjunction with 
% IEEEtran.cls'  captionsoff option. Some IEEE journals/societies require that
% submissions have lists of figures/tables at the end of the paper and that
% figures/tables without any captions are placed on a page by themselves at
% the end of the document. If needed, the draftcls IEEEtran class option or
% \CLASSINPUTbaselinestretch interface can be used to increase the line
% spacing as well. Be sure and use the nomarkers option of endfloat to
% prevent endfloat from "marking" where the figures would have been placed
% in the text. The two hack lines of code above are a slight modification of
% that suggested by in the endfloat docs (section 8.4.1) to ensure that
% the full captions always appear in the list of figures/tables - even if
% the user used the short optional argument of \caption[]{}.
% IEEE papers do not typically make use of \caption[]'s optional argument,
% so this should not be an issue. A similar trick can be used to disable
% captions of packages such as subfig.sty that lack options to turn off
% the subcaptions:
% For subfig.sty:
% \let\MYorigsubfloat\subfloat
% \renewcommand{\subfloat}[2][\relax]{\MYorigsubfloat[]{#2}}
% However, the above trick will not work if both optional arguments of
% the \subfloat command are used. Furthermore, there needs to be a
% description of each subfigure *somewhere* and endfloat does not add
% subfigure captions to its list of figures. Thus, the best approach is to
% avoid the use of subfigure captions (many IEEE journals avoid them anyway)
% and instead reference/explain all the subfigures within the main caption.
% The latest version of endfloat.sty and its documentation can obtained at:
% http://www.ctan.org/pkg/endfloat
%
% The IEEEtran \ifCLASSOPTIONcaptionsoff conditional can also be used
% later in the document, say, to conditionally put the References on a 
% page by themselves.




% *** PDF, URL AND HYPERLINK PACKAGES ***
%
%\usepackage{url}
% url.sty was written by Donald Arseneau. It provides better support for
% handling and breaking URLs. url.sty is already installed on most LaTeX
% systems. The latest version and documentation can be obtained at:
% http://www.ctan.org/pkg/url
% Basically, \url{my_url_here}.





% *** Do not adjust lengths that control margins, column widths, etc. ***
% *** Do not use packages that alter fonts (such as pslatex).         ***
% There should be no need to do such things with IEEEtran.cls V1.6 and later.
% (Unless specifically asked to do so by the journal or conference you plan
% to submit to, of course. )

\usepackage{hyperref}

\usepackage{amsmath}
% correct bad hyphenation here
\hyphenation{op-tical net-works semi-conduc-tor}


\begin{document}
%
% paper title
% Titles are generally capitalized except for words such as a, an, and, as,
% at, but, by, for, in, nor, of, on, or, the, to and up, which are usually
% not capitalized unless they are the first or last word of the title.
% Linebreaks \\ can be used within to get better formatting as desired.
% Do not put math or special symbols in the title.
\title{Social Network Graph Representation For Information Operation Detection}
%
%
% author names and IEEE memberships
% note positions of commas and nonbreaking spaces ( ~ ) LaTeX will not break
% a structure at a ~ so this keeps an author's name from being broken across
% two lines.
% use \thanks{} to gain access to the first footnote area
% a separate \thanks must be used for each paragraph as LaTeX2e's \thanks
% was not built to handle multiple paragraphs
%
%
%\IEEEcompsocitemizethanks is a special \thanks that produces the bulleted
% lists the Computer Society journals use for "first footnote" author
% affiliations. Use \IEEEcompsocthanksitem which works much like \item
% for each affiliation group. When not in compsoc mode,
% \IEEEcompsocitemizethanks becomes like \thanks and
% \IEEEcompsocthanksitem becomes a line break with idention. This
% facilitates dual compilation, although admittedly the differences in the
% desired content of \author between the different types of papers makes a
% one-size-fits-all approach a daunting prospect. For instance, compsoc 
% journal papers have the author affiliations above the "Manuscript
% received ..."  text while in non-compsoc journals this is reversed. Sigh.

\author{Michael~Shell,~\IEEEmembership{Member,~IEEE,}
        John~Doe,~\IEEEmembership{Fellow,~OSA,}
        and~Jane~Doe,~\IEEEmembership{Life~Fellow,~IEEE}% <-this % stops a space
\IEEEcompsocitemizethanks{\IEEEcompsocthanksitem M. Shell was with the Department
of Electrical and Computer Engineering, Georgia Institute of Technology, Atlanta,
GA, 30332.\protect\\
% note need leading \protect in front of \\ to get a newline within \thanks as
% \\ is fragile and will error, could use \hfil\break instead.
E-mail: see http://www.michaelshell.org/contact.html
\IEEEcompsocthanksitem J. Doe and J. Doe are with Anonymous University.}% <-this % stops an unwanted space
\thanks{Manuscript received April 19, 2005; revised August 26, 2015.}}

% note the % following the last \IEEEmembership and also \thanks - 
% these prevent an unwanted space from occurring between the last author name
% and the end of the author line. i.e., if you had this:
% 
% \author{....lastname \thanks{...} \thanks{...} }
%                     ^------------^------------^----Do not want these spaces!
%
% a space would be appended to the last name and could cause every name on that
% line to be shifted left slightly. This is one of those "LaTeX things". For
% instance, "\textbf{A} \textbf{B}" will typeset as "A B" not "AB". To get
% "AB" then you have to do: "\textbf{A}\textbf{B}"
% \thanks is no different in this regard, so shield the last } of each \thanks
% that ends a line with a % and do not let a space in before the next \thanks.
% Spaces after \IEEEmembership other than the last one are OK (and needed) as
% you are supposed to have spaces between the names. For what it is worth,
% this is a minor point as most people would not even notice if the said evil
% space somehow managed to creep in.



% The paper headers
\markboth{Journal of \LaTeX\ Class Files,~Vol.~14, No.~8, August~2015}%
{Shell \MakeLowercase{\textit{et al.}}: Bare Demo of IEEEtran.cls for Computer Society Journals}
% The only time the second header will appear is for the odd numbered pages
% after the title page when using the twoside option.
% 
% *** Note that you probably will NOT want to include the author's ***
% *** name in the headers of peer review papers.                   ***
% You can use \ifCLASSOPTIONpeerreview for conditional compilation here if
% you desire.



% The publisher's ID mark at the bottom of the page is less important with
% Computer Society journal papers as those publications place the marks
% outside of the main text columns and, therefore, unlike regular IEEE
% journals, the available text space is not reduced by their presence.
% If you want to put a publisher's ID mark on the page you can do it like
% this:
%\IEEEpubid{0000--0000/00\$00.00~\copyright~2015 IEEE}
% or like this to get the Computer Society new two part style.
%\IEEEpubid{\makebox[\columnwidth]{\hfill 0000--0000/00/\$00.00~\copyright~2015 IEEE}%
%\hspace{\columnsep}\makebox[\columnwidth]{Published by the IEEE Computer Society\hfill}}
% Remember, if you use this you must call \IEEEpubidadjcol in the second
% column for its text to clear the IEEEpubid mark (Computer Society jorunal
% papers don't need this extra clearance.)



% use for special paper notices
%\IEEEspecialpapernotice{(Invited Paper)}



% for Computer Society papers, we must declare the abstract and index terms
% PRIOR to the title within the \IEEEtitleabstractindextext IEEEtran
% command as these need to go into the title area created by \maketitle.
% As a general rule, do not put math, special symbols or citations
% in the abstract or keywords.
\IEEEtitleabstractindextext{%
\begin{abstract}
This work integrates insights from three domains: graph theory, social network analysis and behavioral studies. Building on these perspectives, we articulate the central research goal: developing enhanced graph representations of social networks for io detection. While numerous io detection methods have been proposed, little research has focused on designing graph representations of social networks that are suitable for newer methods. We address the challenges of constructing relational models from raw data, and of validating results against established io detection approaches.

Our findings demonstrate how improved graph representations can support both existing and emerging io detection methods, helping to optimize their effectiveness. This work contributes to the broader fight against information operations and the spread of misinformation.
\end{abstract}

% Note that keywords are not normally used for peerreview papers.
\begin{IEEEkeywords}
Computer Society, IEEE, IEEEtran, journal, \LaTeX, paper, template.
Social Network, Graph Representation, Information Operations, Influence Campaigns,  Coordinated Inauthentic Behavior
\end{IEEEkeywords}}


% make the title area
\maketitle


% To allow for easy dual compilation without having to reenter the
% abstract/keywords data, the \IEEEtitleabstractindextext text will
% not be used in maketitle, but will appear (i.e., to be "transported")
% here as \IEEEdisplaynontitleabstractindextext when the compsoc 
% or transmag modes are not selected <OR> if conference mode is selected 
% - because all conference papers position the abstract like regular
% papers do.
\IEEEdisplaynontitleabstractindextext
% \IEEEdisplaynontitleabstractindextext has no effect when using
% compsoc or transmag under a non-conference mode.



% For peer review papers, you can put extra information on the cover
% page as needed:
% \ifCLASSOPTIONpeerreview
% \begin{center} \bfseries EDICS Category: 3-BBND \end{center}
% \fi
%
% For peerreview papers, this IEEEtran command inserts a page break and
% creates the second title. It will be ignored for other modes.
\IEEEpeerreviewmaketitle






% \IEEEraisesectionheading{\section{Introduction}\label{sec:introduction}}
% \IEEEPARstart{D}{igital} platforms have become a central arena for Information Operations (IO). IO refers to organized efforts by state or non-state actors to manipulate public opinion and distort political or social discourse through digital platforms. These operations commonly employ disinformation, propaganda, and networks of inauthentic accounts to amplify specific narratives and create misleading perceptions of consensus \cite{weedon2017information}.

% A central mechanism enabling the effectiveness of IO is Coordinated Inauthentic Behavior (CIB), defined as groups of accounts working together to mislead users about their identity or intent \cite{gleicher2018coordinated, mannocci2024detection}. Rather than acting independently, coordinated actors perform similar or complementary actions in pursuit of a shared objective while concealing their coordination.

% Such coordinated activity manifests through observable digital traces, commonly referred to as co-actions. These include co-sharing or retweeting the same content, co-commenting or liking the same posts, posting highly similar textual or visual material, mentioning the same accounts, and using identical hashtags or URLs within short time windows \cite{luceri2024unmasking, mannocci2024detection}. These behavioral similarities provide indirect but measurable signals of coordination.

% Recent advances in artificial intelligence, particularly large language models, have further lowered the barriers for conducting IO by enabling large-scale generation of persuasive and human-like content. This development increases both the scale and sophistication of coordinated campaigns, rendering traditional content-based detection methods increasingly insufficient \cite{fredheim2024assessing}. Consequently, detection efforts must increasingly rely on relational and structural patterns of interaction rather than solely on textual features.

% Beyond their immediate effects on public discourse, IO pose significant long-term risks to democratic institutions. By eroding trust in political processes and amplifying social polarization, sustained influence campaigns can reshape public perception over extended periods. A prominent example is the \textit{Stop the Steal} campaign surrounding the 2020 U.S. Presidential Election, in which coordinated actors systematically amplified false claims of voter fraud across multiple platforms. The campaign progressed through distinct stages, from seeding conspiracy theories and broadcasting narratives to convergence around a common message and ultimately offline mobilization, culminating in the January 6, 2021 Capitol attack \cite{coordinated2024Capitol, LongFuse2021USElections}. Notably, even a year after the election, a substantial portion of the U.S. population continued to doubt the legitimacy of its outcome \cite{stopthesteal2022USelections}. This case underscores how coordinated digital activity can translate into real-world consequences, and highlights the critical importance of detecting IO early and reliably.

% Graph-based representations of social networks have emerged as a promising approach for capturing the structural and relational signals of coordination. By modeling users as nodes and their interactions or behavioral similarities as edges, such representations enable detection methods to exploit the topology of coordinated activity rather than relying solely on content features. However, despite the growing body of work on IO detection algorithms, surprisingly little attention has been devoted to the construction of the underlying graph itself. Most existing approaches treat the social network graph as a fixed, predefined input, without systematically examining how different representational choices, such as edge definitions, weighting schemes, or graph compression strategies, affect detection performance. This gap is consequential: the quality and structure of the graph representation directly determine what signals are available to any downstream detection method.

% This work addresses that gap. Rather than proposing a new detection algorithm or end-to-end pipeline, we focus on the fundamental problem of graph construction from raw social media data, and examine how alternative representations influence the effectiveness of known IO detection approaches. Specifically, this paper is guided by the following research questions:

% \begin{itemize}
%     \item \textbf{RQ1:} How does the choice of social network graph representation affect the accuracy of different detection methods?
%     \item \textbf{RQ2:} Which combination of graph representation and IO detection method achieves the best performance?
%     \item \textbf{RQ3:} Can graph representations be used with graph learning techniques to provide a simple and effective framework for IO detection?
% \end{itemize}



\IEEEPARstart{D}{igital} platforms have become a central arena for Information Operations (IO), organized efforts by state or non-state actors to manipulate public opinion and distort political or social discourse through digital platforms. These operations commonly employ disinformation, propaganda, and networks of inauthentic accounts to amplify specific narratives and create misleading perceptions of consensus \cite{weedon2017information, troll2020detection}. A central mechanism enabling their effectiveness is Coordinated Inauthentic Behavior (CIB), defined as groups of accounts working together to mislead users about their identity or intent \cite{gleicher2018coordinated, mannocci2024detection}. Rather than acting independently, coordinated actors perform similar or complementary actions in pursuit of a shared objective while concealing their coordination.

Beyond their immediate effects on public discourse, IO pose significant long-term risks to democratic institutions. By eroding trust in political processes and amplifying social polarization, sustained influence campaigns can reshape public perception over extended periods. A prominent example is the \textit{Stop the Steal} campaign surrounding the 2020 U.S. Presidential Election, in which coordinated actors systematically amplified false claims of voter fraud across multiple platforms. The campaign progressed through distinct stages, from seeding conspiracy theories and broadcasting narratives to convergence around a common message and ultimately offline mobilization, culminating in the January 6, 2021 Capitol attack \cite{coordinated2024Capitol, LongFuse2021USElections}. Furthermore, even a year after the election, a substantial portion of the U.S. population continued to doubt the legitimacy of its outcome \cite{stopthesteal2022USelections}.

Across such campaigns, coordinated activity manifests through observable digital traces, commonly referred to as co-actions. These include co-sharing or retweeting the same content, co-commenting or liking the same posts, posting highly similar textual or visual material, mentioning the same accounts, and using identical hashtags or URLs within short time windows \cite{luceri2024unmasking, mannocci2024detection}. These behavioral similarities provide indirect but measurable signals of coordination.

Recent advances in artificial intelligence, particularly large language models, have further lowered the barriers for conducting IO by enabling large-scale generation of persuasive and human-like content \cite{burghardt2024large}. This development increases both the scale and sophistication of coordinated campaigns, rendering traditional content-based detection methods increasingly insufficient \cite{fredheim2024assessing, luceri2024leveraging}. Consequently, detection efforts must increasingly rely on relational and structural patterns of interaction rather than solely on textual features.


Graph-based representations of social networks have emerged as a promising approach for capturing the structural and relational signals of coordination \cite{minici2025iohunter, gabriel2023inductive, pacheco2021uncovering, smith2021automatic}. By modeling users as nodes and their interactions or behavioral similarities as edges, such representations enable detection methods to exploit the topology of coordinated activity rather than relying solely on content features. However, most existing approaches construct this graph from raw follow or repost interactions over the entire user base, without distinguishing influential actors from peripheral participants, and without incorporating the narrative context that characterizes specific campaigns. Structural signals of coordination are thus diluted by the volume of low-signal nodes. As a result, few studies have examined how core design decisions, such as node selection, edge weighting, and topic-level scoping, affect detection performance \cite{coscia2019impact, backbone2022collectiveBehavior, graphCompression2018survey}.

This work addresses that gap by introducing the \textit{Backbone Graph}: a compact, narrative-aware representation that connects only the most influential authors within each topic through their behavioral similarities over shared entities (hashtags, URLs, account mentions, and named entities), weighted by TF-IDF. Rather than proposing a new detection algorithm, we examine whether this representation alone exposes the structural position of IO drivers. Specifically, we test the following hypotheses:

\begin{itemize}
    \item \textbf{H1 (Dense Community):} IO drivers are more likely to appear in densely connected regions of the Backbone Graph.
    \item \textbf{H2 (Community Heterogeneity):} IO drivers are non-uniformly distributed across communities, concentrating in specific narrative-based groups.
    \item \textbf{H3 (Matching Cohort):} The key-author score identifies IO drivers beyond what is explained by general activity and structural features (e.g., post count, degree centrality).
    \item \textbf{H3.1 (Narrative Amplification):} IO drivers tend to lead narratives within their topics, reflected in higher key-author scores than non-IO users.
\end{itemize}



\section{Related Work}

\subsection{Information Operations}

Coordinated information operations (IO) have been studied as a distinct phenomenon since at least 2017, when Facebook~\cite{weedon2017information} 
defined them as coordinated actions intended to distort political sentiment, distinguished from organic misinformation by intent, coordination, and 
inauthentic behavior. The U.S. Office of the Director of National Intelligence~\cite{activities2017intentions} documented state-sponsored campaigns in parallel, and these two reports established the vocabulary that detection research has used ever since. More recent
work~\cite{fredheim2025assessing,zaman2025social} examines how LLMs reshape the threat landscape, but the underlying structure of IO, coordinated actors
amplifying shared narratives, has remained stable. Empirical studies of specific campaigns, such as IO on Twitter around major elections~\cite{coordinated2019Twitter}, have helped ground these definitions in practice, and publicly released labeled datasets of state-linked campaigns now enable controlled comparisons across detection methods~\cite{seckin2025labeled}.


\subsection{Content-Based Detection}


Many detection systems operate on what accounts say. TrollHunter2020~\cite{jachim2021trollhunter2020} uses verb and noun co-occurrence patterns to spot emerging trolling narratives on Twitter during the 2020 U.S. elections, though its reliance on surface lexical patterns breaks down when adversaries paraphrase. Sociolinguistic features, such as stylistic and interaction-behavior differences, have also been examined as signals of inauthenticity~\cite{SocioLinguistic2024coordinated}, and LLMs have more recently been used to extract narrative frames directly from post content~\cite{burghardt2024large}.

What these approaches share is a focus on the content that accounts produce, rather than on how accounts are structurally organized within the broader interaction network.

\subsection{Relational and Graph-Based Approaches}

Graph-based methods fill this gap by modeling the relationships between accounts. Recent survey document the rapid growth of this literature and its main design patterns~\cite{mannocci2024detection}. Foundational work in this space established the analytical vocabulary that later systems built on: Giglietto et al.~\cite{giglietto2020village} introduced coordinated link sharing behavior as a detectable signature of inauthentic amplification, Nizzoli et al.~\cite{coordinated2021UK} operationalized coordination detection around shared actions within narrow time windows, and Magelinski et al.~\cite{magelinski2022synchronized} proposed a synchronized-action framework that formalized what counts as a detectable coordination event. Building on these ideas, Pacheco et al.~\cite{pacheco2021uncovering} introduced an
unsupervised framework that builds bipartite user-entity graphs from behavioral traces such as co-retweets, URL sharing, and hashtag usage, and projects them
into weighted similarity networks where edges above a Jaccard threshold reveal coordination clusters. The framework works, but the reliance on
hand-tuned thresholds raises questions about robustness and reproducibility. A related line of work frames detection at the structural level: coordinated accounts tend to cluster both in the dense core of the co-action graph and within specific communities, and community-detection techniques have been used to surface these clusters directly~\cite{gang2020coordinating}. This observation motivates the use of placement-based features, such as k-core membership and community assignment, as diagnostic signals of coordinated activity.

Luceri et al.~\cite{luceri2024unmasking} sidestep this issue by replacing edge filtering with node pruning based on eigenvector centrality over a 
fused similarity network, reporting strong performance across several datasets. Gabriel et al.~\cite{gabriel2023inductive} go further, showing 
that structural representations such as node2vec, PageRank, and HITS outperform content-based features in inductive graph learning. This is evidence
that much of the signal lives in the relational structure alone.

A recent variant combines graph structure with LLMs: Luceri et al.~\cite{luceri2024leveraging} translate retweet-graph topology into textual descriptions enriched with centrality and interaction metadata, feeding them to open-source LLMs for campaign classification. Although effective, the translation from graph to text inevitably discards structural information, motivating approaches that operate on the graph directly.

Hybrid approaches that combine structure and content have also done well.
Smith et al.~\cite{smith2021automatic} merge NLP, graph analytics, and causal inference 
to study coordinated activity in the 2017 French presidential election, 
achieving high precision and surfacing influential accounts that were later confirmed by external reports. IOHunter~\cite{minici2025iohunter} 
pairs sentence-BERT embeddings with a cross-attention GNN, improving F1 by up to 20\% over prior baselines and generalizing well across domains. 
A common ingredient in these methods is similarity graphs built from shared entities, where features capture pairwise resemblance between
accounts: co-URL, co-hashtag, or co-retweet patterns encoded via TF-IDF or learned embeddings. Open-source toolkits have also consolidated such techniques into reusable detection pipelines~\cite{toolkit2024coordination}. These methods, however, build similarity graphs over the full user base and do not restrict the graph to a curated set of influential actors. The Backbone Graph introduced below differs in precisely this choice, operating on a compact set of topic-specific key authors rather than on all interacting users.


\section{Methodology}
\subsection{Datasets}
We use publicly available datasets of state-linked information operations introduced by Seckin et al.~\cite{seckin2025labeled} . The datasets include coordinated campaigns identified by social media platforms, along with the activity timelines of accounts associated with each campaign. In this work, we examine several campaigns from this dataset, including the Ecuador campaign, the [CAMPAIGN X] campaign, and the [CAMPAIGN Y] campaign. Each campaign dataset also includes a control group of organic accounts that discussed similar topics during the same time period, allowing us to compare coordinated and non-coordinated activity.

\subsection{Topic Extraction}
To identify the main topics within each campaign, we applied BERTopic~\cite{grootendorst2022bertopic} separately to each dataset. Since the optimal level of topic granularity is not known in advance, we experimented with multiple values of the min\_topic\_size parameter, which controls the minimum number of documents required to form a topic, and selected the value that yielded the most coherent and meaningful topic structure for each corpus. The resulting topics serve as the basis for identifying key entities, posts, and authors within each discourse.
\subsection{Key Entity Identification}
For each topic, we extract the entities appearing in its posts, including named entities, hashtags, and mentions. To measure the relative importance of each entity within a topic, we compute a TF-IDF based score.
Formally, let $e$ denote an entity and $t$ a topic. Each topic is represented as a concatenated document of all entities mentioned in its posts. The TF-IDF score is then computed as:
\begin{equation}
    \text{TF-IDF}_{e,t} = \text{TF}_{e,t} \times \text{IDF}_{e}
\end{equation}

where $\text{TF}_{e,t}$ is the frequency of entity $e$ within the topic document of $t$, and $\text{IDF}_{e}$ is the inverse document frequency of $e$ across all topics:

\begin{equation}
    \text{IDF}_{e} = \log\left(\frac{|\text{all\_topics}|}{|\text{topics containing } e|}\right)
\end{equation}

This score assigns higher weight to entities that are prominent within a specific topic and relatively rare across the corpus.
\section{Conclusion}
The conclusion goes here.





% if have a single appendix:
%\appendix[Proof of the Zonklar Equations]
% or
%\appendix  % for no appendix heading
% do not use \section anymore after \appendix, only \section*
% is possibly needed

% use appendices with more than one appendix
% then use \section to start each appendix
% you must declare a \section before using any
% \subsection or using \label (\appendices by itself
% starts a section numbered zero.)
%


\appendices
\section{Proof of the First Zonklar Equation}
Appendix one text goes here.

% you can choose not to have a title for an appendix
% if you want by leaving the argument blank
\section{}
Appendix two text goes here.


% use section* for acknowledgment
\ifCLASSOPTIONcompsoc
  % The Computer Society usually uses the plural form
  \section*{Acknowledgments}
\else
  % regular IEEE prefers the singular form
  \section*{Acknowledgment}
\fi


The authors would like to thank...


% Can use something like this to put references on a page
% by themselves when using endfloat and the captionsoff option.
\ifCLASSOPTIONcaptionsoff
  \newpage
\fi



% trigger a \newpage just before the given reference
% number - used to balance the columns on the last page
% adjust value as needed - may need to be readjusted if
% the document is modified later
%\IEEEtriggeratref{8}
% The "triggered" command can be changed if desired:
%\IEEEtriggercmd{\enlargethispage{-5in}}

% references section

% can use a bibliography generated by BibTeX as a .bbl file
% BibTeX documentation can be easily obtained at:
% http://mirror.ctan.org/biblio/bibtex/contrib/doc/
% The IEEEtran BibTeX style support page is at:
% http://www.michaelshell.org/tex/ieeetran/bibtex/
%\bibliographystyle{IEEEtran}
% argument is your BibTeX string definitions and bibliography database(s)
%\bibliography{IEEEabrv,../bib/paper}
%
% <OR> manually copy in the resultant .bbl file
% set second argument of \begin to the number of references
% (used to reserve space for the reference number labels box)
\begin{thebibliography}{1}

\bibitem{IEEEhowto:kopka}
H.~Kopka and P.~W. Daly, \emph{A Guide to \LaTeX}, 3rd~ed.\hskip 1em plus
  0.5em minus 0.4em\relax Harlow, England: Addison-Wesley, 1999.

\end{thebibliography}

% biography section
% 
% If you have an EPS/PDF photo (graphicx package needed) extra braces are
% needed around the contents of the optional argument to biography to prevent
% the LaTeX parser from getting confused when it sees the complicated
% \includegraphics command within an optional argument. (You could create
% your own custom macro containing the \includegraphics command to make things
% simpler here.)
%\begin{IEEEbiography}[{\includegraphics[width=1in,height=1.25in,clip,keepaspectratio]{mshell}}]{Michael Shell}
% or if you just want to reserve a space for a photo:

\begin{IEEEbiography}{Michael Shell}
Biography text here.
\end{IEEEbiography}

% if you will not have a photo at all:
\begin{IEEEbiographynophoto}{John Doe}
Biography text here.
\end{IEEEbiographynophoto}

% insert where needed to balance the two columns on the last page with
% biographies
%\newpage

\begin{IEEEbiographynophoto}{Jane Doe}
Biography text here.
\end{IEEEbiographynophoto}

% You can push biographies down or up by placing
% a \vfill before or after them. The appropriate
% use of \vfill depends on what kind of text is
% on the last page and whether or not the columns
% are being equalized.

%\vfill

% Can be used to pull up biographies so that the bottom of the last one
% is flush with the other column.
%\enlargethispage{-5in}
\bibliographystyle{unsrt}
\bibliography{references}



% that's all folks
\end{document}


